\documentclass[10pt]{article}
\usepackage[T1]{fontenc}
\usepackage{lmodern,amsmath,amssymb,graphicx}
\ifdefined\pdfstartlink
  \newcommand{\linkto}[2]{\leavevmode\pdfstartlink attr {/Border [0 0 0]} user {/Subtype /Link /A << /S /URI /URI (\pdfescapestring{\detokenize{#1}}) >>}#2\pdfendlink}
  \newcommand{\marksection}[2]{\pdfdest name {#1} xyz\pdfoutline goto name {#1} {#2}}
  \pdfinfo{/Title (CrytalDrift: One-step generation in the geometry of periodic crystals) /Author (AmirM.MohammadHosseini) /Subject (Invariant training fields, periodic geometry and Wyckoff-conditioned crystal generation)}
\else
  \usepackage[hidelinks]{hyperref}
  \hypersetup{pdftitle={CrytalDrift: One-step generation in the geometry of periodic crystals},pdfauthor={AmirM.MohammadHosseini}}
  \newcommand{\linkto}[2]{\href{#1}{#2}}\newcommand{\marksection}[2]{\pdfbookmark[1]{#2}{#1}}
\fi
\setlength{\textwidth}{6.43in}
\setlength{\oddsidemargin}{0.035in}
\setlength{\textheight}{9.12in}
\setlength{\topmargin}{-.53in}
\setlength{\parindent}{1.05em}
\setlength{\emergencystretch}{1.5em}
\setlength{\textfloatsep}{11pt plus 2pt minus 2pt}
\renewcommand{\topfraction}{.92}
\renewcommand{\bottomfraction}{.8}
\renewcommand{\textfraction}{.08}
\renewcommand{\floatpagefraction}{.76}
\setcounter{topnumber}{3}\setcounter{totalnumber}{4}
\AddToHook{env/figure/before}{\par}
\AddToHook{env/figure/after}{\par}
\AddToHook{env/table/before}{\par}
\AddToHook{env/table/after}{\par}
\title{\vspace{-1.6em}\bfseries CrytalDrift\\[2pt]\large One-step generation in the geometry of periodic crystals}
\author{AmirM.MohammadHosseini\\\small Department of Biotechnology, University of Tehran}
\date{}
\begin{document}
\maketitle
\vspace{-1.4em}
\begin{abstract}
Can one neural evaluation recover a periodic crystal when each training example can be written in a different lattice basis? CrytalDrift jointly updates the cell and fractional sites of a random seed with a periodic graph network. During training, attraction to reference structures and repulsion from independent model draws act on species-resolved radial shells and local angular power spectra. We prove that each query and reference may independently change its unimodular basis, Cartesian frame, origin and site order while leaving the empirical training field unchanged. On all 9,046 MP-20 test structures, the angular generator recovered 31.79\% of targets in twenty raw draws, averaged over three fits. The same network trained with pair shells recovered 31.30\%, the earlier fixed-basis generator 32.91\%, and released one-step uFlowCSP 79.32\%. Requiring native volume agreement reduced angular and radial recovery to 21.24\% and 24.60\%, respectively. In a balanced timing replay at 100 ms per target, the angular model recovered 32 of 128 targets, compared with 13 for a matched 16-call flow and 39 for its zero-call seed prior. Supplying the occupied Wyckoff sites changes the task. One template-conditioned proposal recovered 50.29\% of targets, compared with 45.90\% from its seed prior alone. Under a train-fitted composition prior, 95.59\% of 10,000 unfiltered angular proposals passed the structural validity test. None of the first sixteen passed every criterion of a prespecified relaxation and competing-phase screen. These results quantify what representation-independent drifting preserves and what it costs in predictive and physical tests.

\end{abstract}

\section{Learning from different descriptions of the same crystal}
A cubic cell can be redrawn as a skew parallelepiped, changing both its lattice coordinates and every fractional site. A training objective that compares these arrays directly can exert a force between two descriptions of the same material. With a bank of training and generated samples, each member can have its own choice of basis, origin and frame. Equivariance of a network acting on a single input leaves this comparison problem open.

\begin{figure}[!htbp]\centering
\includegraphics[width=\linewidth]{csp-interleaved.pdf}
\caption{What one call buys at fixed sample count and fixed elapsed time. Each method receives the same 128 MP-20 cell compositions and a stream capped at 20 completed proposals per target. For each target we vary the order in which methods run. The physical networks share seed draws. The right panel counts only proposals already returned to host memory by the deadline. Matching happens afterward. Shading gives a pointwise 95\% target-bootstrap interval for the angular model.}
\label{fig:v2balanced}\end{figure}

We adapt the training-time attraction and repulsion of Deng et al.\ to periodic atomic environments \cite{deng}. In CrytalDrift, a scalar affinity determines how much each reference contributes. Differences of physical features determine the direction of the update. A separate graph network takes one seed and proposes a cell and all of its sites in one call. When an occupied Wyckoff template is supplied, a crystallographic construction first draws a compatible seed. These three operations answer different questions: which crystals are similar, how a proposal moves, and which symmetry operations it must retain. Figure~\ref{fig:method} follows them from inputs to outputs.
\begin{figure}[!htbp]\centering\includegraphics[width=\linewidth]{method.pdf}
\caption{The formula-only experiment (A) supplies the elements and their cell multiplicities. The template experiment (C) also supplies a group in a specified setting and its species-labelled occupied orbits. The upper path returns a lattice and fractional coordinates after one neural evaluation. The two banks and the feature-space drift below enter training only.}
\label{fig:method}\end{figure}

For molecules, SymDrift compares aligned coordinates or embeds each structure in its sorted lists of element-pair distances \cite{symdrift}. Our invariant comparison must also account for the translation lattice, its nonzero self-images and the incidence of several neighbors at one center. Local power spectra and whole-structure pooling have a long history in atomistic modeling \cite{soap,whole,degenerate}. LEGO-xtal uses analogous harmonic contractions during pre-relaxation toward specified local environments \cite{lego}. Here the contractions weight and direct a distributional training update on jointly proposed cells and sites. Kernel-gradient drifting obtains directions by differentiating a kernel and has identification results under its own assumptions \cite{kernelgrad}. TokenDrift refines text generation in an embedding and Drift-React generates endpoint-conditioned molecular paths \cite{tokendrift,driftreact}.

\marksection{theory}{Periodic training field and geometric tests}
\section{A field on physical periodic environments}\label{sec:physical-theory}
\paragraph{The finite training update.}
Let $q_i=\Phi(G_\theta(\xi_i,c))$ denote a query feature at exact cell composition $c$.
In each of eight composition blocks, four independent empirical draws with replacement
supply features $p_j$, and four independent model seeds supply features $n_j$.
Write $y=(p,n)$, $H_{ij}=-\|q_i-y_j\|_2/\tau$, and
\[
 a_{ij}=\frac{\exp H_{ij}}{\sqrt{\sum_u\exp H_{uj}\sum_v\exp H_{iv}}}.
\]
The sums run over the four queries and the eight reference columns of this block.
Split $a=(a^+,a^-)$ along the reference axis and write $m_i^\pm=\sum_j a^\pm_{ij}$.
The feature displacement and regression target are
\begin{align}
 D_i^\tau&=m_i^-\sum_j a^+_{ij}p_j-m_i^+\sum_j a^-_{ij}n_j,\qquad
 V_i=\tfrac12\sum_{\tau\in\{0.3,0.8\}}
 \frac{D_i^\tau}{\sqrt{\langle\|D^\tau\|^2\rangle+10^{-8}}},\label{eq:v2field}\\
 {\cal L}&=\tfrac12\left\langle
 \big\|q_i-\operatorname{stop}(q_i+V_i)\big\|^2\right\rangle .\label{eq:v2loss}
\end{align}
The mean in each denominator covers every query in the entire minibatch.
Equation~\eqref{eq:v2loss} gives parameter gradient $-\langle J_i^\top V_i\rangle$
with $J_i=\partial_\theta\Phi(G_\theta(\xi_i,c))$. We cache positive features and detach
negative draws and target construction. The Laplacian affinity
sets the weights. Reference feature differences supply the directions. A diffusion
transition would instead specify how to sample a state given a previous state.
At inference the network never visits the training banks.

\paragraph{Why the representation enters the field.}
Consider the row basis $L=I$ and a shear $U=I+e_1e_2^\top$ of the same one-site
unit crystal at the origin. A feature containing the nine entries of $L$ assigns unit
distance between these descriptions. In a one-query, one-positive, one-negative example
all three may describe this same crystal. Replacing only the positive basis by $U$
changes the unnormalized displacement from zero to
$a^+a^-\operatorname{vec}(U-I)$. Equivariance of a generator acting on the query alone
does not correct this dependence on representations chosen independently for the references.

\paragraph{Radial and angular features.}
For row fractional coordinates $f_i$, cell $L$ in \AA, and species $z_i$, take each image
$d=(f_j-f_i+n)L$ of length below $r_c=4$~\AA. Only $i=j,\ n=0$ is omitted.
If $x=f_j-f_i$, Cauchy--Schwarz gives
$|x_k+n_k|\le r_c\|(L^{-1})_{:k}\|$ for every accepted image.
Enumerating the resulting integer intervals and testing the Euclidean radius therefore
includes all physical neighbors under a shear. Put $s=\|d/r_c\|^2$ and
\[
 b_h(s)=\exp[1-(1-s)^{-1}]\,\exp[-\tfrac12(s-h)^2],
 \qquad 0\leq s<1,\quad h\in\{0,1\},
\]
with zero extension past the cutoff. Use real regular solid harmonics
$T_{\ell m}(d/r_c)$ for $\ell=0,1,2$, with the normalization
\[
 \sum_m T_{\ell m}(u)T_{\ell m}(v)=(\|u\|\|v\|)^\ell
 P_\ell(\widehat u\cdot\widehat v).
\]
The polynomial form at the origin assigns value one for order zero and zero for
positive orders, including at a coincident distinct atom. For each center form
\[
 c_{iah\ell m}=\sum_{e:i\to a} b_h(s_e)\, T_{\ell m}(d_e/r_c),\qquad
 p_{i,ah,bh',\ell}=\sum_m c_{iah\ell m}c_{ibh'\ell m}.
\]
This is a discrete density expansion with radial factor $s^{\ell/2}b_h(s)$ in the unit-vector
harmonic convention. It does not first convolute atoms with Gaussian spatial densities.
Average the upper-triangular products over centers of each element \emph{after} contraction.
Seven ordered neighbor-species ranks and two radial channels give $7\cdot105$ entries at
each order. Ranks refer to increasing atomic number and are compared only within an exact
composition. Missing ranks are padded with zeros. The shared 177-entry control contains
log absolute cell volume, 118 global species counts, the nonzero lattice-image shells,
and the 28 unordered distinct-site species-pair shells divided by twice the site count.
The full descriptor has 2,382 entries. For each block subtract its training centroid and
divide by its training root-mean-square Euclidean deviation, floored at $0.01$.
The common blocks retain identical scales under ablation.

\paragraph{Independent invariance of every query and reference.}
Let $U\in GL(3,\mathbb Z)$ change the lattice basis, $Q\in O(3)$ change the Cartesian
frame and $\pi$ permute the sites with their species. Put
$L'=ULQ$, $f'_{\pi(i)}=f_iU^{-1}+t-h_i$ with integer wrapping vectors $h_i$.
The image label $n'=nU^{-1}+h_j-h_i$ is integral and runs bijectively over
$\mathbb Z^3$. Its transformed edge is $d'=dQ$. It preserves radii, species and
the nontrivial-self-edge rule. Every harmonic coefficient at order $\ell$ therefore
transforms by the same orthogonal real representation on its $m$ index. Reflection
contributes parity $(-1)^\ell$ upon decomposing the orthogonal frame change into inversion
and a proper rotation. Contraction at the same order cancels that action, including parity.
Pooling commutes with $\pi$. The edge bijection preserves each radial shell, all centers
have the same lattice-image shell, and $|\det U|=|\det Q|=1$ preserves the log volume.
The standardized feature vector is unchanged. Applying this argument to each crystal
with its \emph{own} transformations leaves every affinity and both sides of
Eq.~\eqref{eq:v2field} unchanged. Enlarging the cell by a nonunimodular supercell map
is outside the assertion.

\paragraph{Equivariance of the one-call decoder.}
The network embeds elements, aggregates scalar messages from edge lengths, and passes
the site embedding, aggregate and log specific volume to its hidden site head. Its radial
input uses Gaussian channels in $s$ and its message weight is
$w_{ij}=(1-s)^3/(1+\sum_{e:i}(1-s_e)^3)$ for edges inside the cutoff.
Bounded scalar edge coefficients multiply the physical vectors to give uncentered
site shifts $u_i$. Write $v_i=u_i-\sum_j u_j/N$ over the $N$ occupied sites. A second
bounded coefficient multiplies $(d/r_c)^\top(d/r_c)$ and the result is summed with the
same edge weights and divided by $N$. Denote this symmetric tensor by $B$. A pooled
scalar $\alpha\in[-1,1]$ gives $E=(\alpha I+B)/4$. The decoder returns
\[
 L_{\rm out}=L\exp E,\qquad
 F_{\rm out}=(F+vL^{-1})\bmod1 .
\]
The fractional displacement uses the \emph{input} lattice before the affine cell update.
Invariant hidden scalars give $v'=PvQ$ and $E'=Q^\top E Q$. Conjugating the exponential
gives $L'_{\rm out}=UL_{\rm out}Q$. Since
$v'(L')^{-1}=PvL^{-1}U^{-1}$, the output sites transform by the same permutation,
rebasing and origin shift. This pathwise statement pairs a given seed with its image.
Pushforward sampling laws agree when the transformed seed law is itself that pushforward.
The study draws physical cells in a specified chart. It does not assume a uniform law
over integer bases. Differentiation with respect to fractional coordinates yields a
covector $g_f$. The row-form Cartesian gradient $g_fL^{-\top}$ corresponds to a
fractional gradient velocity $g_f(LL^\top)^{-1}$ with the inverse metric.
Here the learned displacement is already a physical vector. For scoring and storage the
proposal service converts the output Gram matrix to six log-metric entries and decodes
a canonical row lattice. This last change of Cartesian representative preserves cell
lengths, angles and all physical position comparisons up to a frame change. Literal
frame equivariance of matrix entries refers to the decoder before canonicalization.

\paragraph{Which group is guaranteed.}
For a species-labelled seed symmetry with fractional action $f\mapsto fR+t$,
$RGR^\top=G$ for $G=LL^\top$. Thus $Q=L^{-1}RL$ is orthogonal.
Use the intervening permutation of sites, $U=R^{-1}$, and this $Q$ in the decoder
identity. The reexpressed seed coincides with the initial one. Hence
$RL_{\rm out}=L_{\rm out}Q$, and its site action maps the labelled output set to
itself modulo integers. Each named seed operation survives this map in exact arithmetic.
Some outputs can gain operations or admit coincident sites. Ordinary continuous random
seeds acquire no prescribed nontrivial stabilizer from the argument.

\paragraph{Equilibrium and its scope.}
With fixed model, labels and queries, suppose each positive bank and independent
equal-size negative bank are exchangeable as a pair, jointly over composition blocks.
Exchanging the banks permutes the columns of $a$, interchanges their masses, and reverses
every $D^\tau$ simultaneously. Its global RMS is unchanged, so $\mathbb E[V_i]=0$,
where an individual finite sample can give nonzero drift. The normalized field is
continuous in finite-dimensional input features and each query norm is bounded by the
square root of the total query count. Consequently weak convergence of the conditional
feature laws, for independent banks of fixed sizes, also gives convergence of the
expected field by convergence of finite products and bounded continuous integration.
Neither result supplies a converse or a training convergence rate. Finite cutoffs can
hide all changes to separated sites, and even-parity contractions identify mirror
environments. Vanishing expected drift cannot identify a unique law of physical crystals.

The vector and dyadic updates draw on equivariant graph and periodic deformation models \cite{egnn,gems,vector}. CrystalFlow, FlowMM and DiffCSP model continuous periodic geometry with integration or denoising \cite{crystalflow,flowmm,diffcsp}. One-step generation itself has clear precedents in uFlowCSP and OMatG-flash \cite{uflow,omatg}. The argument above concerns independent representations inside an empirical field as well as the forward map of the generator. We test its practical value against a fixed-basis release, held-out structures, and clocks that include construction of the neighbor graph.
\section{What the field distinguishes}\label{sec:v2checks}
\begin{figure}[!htbp]\centering
\includegraphics[width=\linewidth]{equivalence.pdf}
\caption{A shear changes the drawn unit cell while preserving its periodic lattice.
In the two four-site motifs the six unordered pairwise distances agree exactly.
The two shortest edges have different incidence. The power spectra are computed on
the physical four-site cells using the listed motif coordinates. These are deterministic
mathematical examples.}\label{fig:v2-witness}
\end{figure}
Place identical atoms at vertices $\{0,1,4,6\}$ or $\{0,1,3,7\}$ of a regular
12-gon of radius 1~\AA\ in a 16~\AA\ cubic cell. Set $r_c=2.25$~\AA.
The pairwise separations are the six chords $2\sin(k\pi/12)$, $k=1,\ldots,6$,
once each in either motif. Only the second motif has its two shortest edges incident
to a common vertex. Any distance-preserving labelling must preserve that incidence,
so the motifs are inequivalent. Their radial blocks agree to $1.6\times10^{-16}$.
The unscaled Euclidean differences for orders zero, one and two are respectively
$0.31544$, $0.06769$ and $0.00617$. Even order zero detects the difference because
multiplication before pooling preserves correlations between shells incident to
the same center. Directional harmonics are sufficient for this example but unnecessary
for its separation. In the fixed-basis release's decorated eight-site witness, we
used its unit cell, coordinates and species with $r_c=4$~\AA. Radial shells agreed
within $8.6\times10^{-13}$ and the three power blocks differed by $0.00246$,
$0.01381$ and $0.00647$. A model can distinguish these examples in feature space
without learning how to place their atoms.

Two simple degeneracies survive the full descriptor. In a 24~\AA\ cubic cell, a
two-species motif with sites 6~\AA\ apart and one with sites 9~\AA\ apart have empty
neighborhoods at the 3.4~\AA\ test cutoff. Every block agrees. A four-element
tetrahedron and its reflected image have opposite signed volumes under a fixed ordering
of the species and identical descriptors. Their six labelled edges fit inside a region
smaller than half a lattice period. No wrap ambiguity can restore the sign by a proper
rigid motion. These examples rule out injectivity.

\paragraph{Numerical checks.}
We drew independent integer shears, determinant signs, orthogonal frames, origins,
wraps and site orders for a batch comprising three differentiable query seeds,
three positive references and three independent negative seeds. With FP64 arithmetic
the maximum changes in the feature vector, field and parameter gradient were respectively
$2.2\times10^{-14}$, $1.6\times10^{-14}$ and $5.3\times10^{-14}$.
The maximum physical site discrepancy was $1.3\times10^{-14}$~\AA.
The finite-batch field had norm $1.69$ and swapping the banks reversed it to
$1.9\times10^{-16}$ precision. The integer-edge code agreed with direct exhaustive
enumeration in the unsheared physical cell to $1.4\times10^{-14}$~\AA\ after transformation.
A second implementation contracted Legendre polynomials directly over all ordered
pairs of edges incident to a center. Its maximum difference from the harmonic calculation
was $2.3\times10^{-15}$. A directional derivative of the descriptor differed from a
central finite difference by $6.9\times10^{-11}$. Tests include nonzero self-images,
distinct coincident atoms, interior padding, mixed sizes, a singleton, empty neighborhoods,
smooth cutoff crossing, single-record positive caching, and controls changing one length,
angle or species. The pinned CUDA FP32 run with TF32 disabled differed from the CPU by
at most $1.5\times10^{-6}$ in features and $2.7\times10^{-6}$ in parameter gradients.
We restored all nine validation-selected periodic checkpoints across the three feature
arms and three seeds. For each, 96 independently transformed synthetic input crystals
were arranged into the actual eight four-query training blocks with that arm's fitted
scales and descriptor. The field, parameter gradient, unwrapped Cartesian shift, strain
and transformed outputs agreed to $4.1\times10^{-14}$ in FP64.

\paragraph{Reach of one bounded update.}
The normalized edge weights sum to less than one at each center. Since
$\|d/r_c\|<1$, the sum of the singular values of $B$ is at most one.
Thus $|\log(V_{\rm out}/V)|=|\operatorname{tr}E|\le1$ and the change in log
singular-value aspect ratio of the cell is at most $1/4$. For site $i$ the uncentered
move is bounded by $b_i=\sum_{e:i}w_e\|d_e\|$. A relative move of two sites in the
seed cell is bounded by $b_i+b_j$. On 2,048 random MP-20 \emph{training} targets, we
compared one independent shape/volume/site prior draw to its supplied native target.
The original prototype seed exceeded the necessary volume and aspect bounds on 120
and 1,112 cases. For the fitted physical prior these counts were 1 and 1,276.
For the 546 cells with at least two uniquely labelled species, a lower bound on the
required relative shift exceeded $b_i+b_j$ on 108 original seeds and 78 fitted seeds.
We used the wrapped fractional displacement multiplied by the smallest seed-cell
singular value as the lower bound. These tests concern a single fixed seed and native
representative per target. The evaluation matcher can rescale and rebase, and another
draw can start nearer a target.

In a separate validation sensitivity, we multiplied only the traceless part of $E$
by four. With $A=L^{-1}L_{\rm out}$ the widened output is
$L_{\rm out}(\det(A)^{-1/3}A)^3$, which keeps the volume change and site update of the
same parameters. Conjugation still commutes with this map. At 256 updates, the wide
minus original validation geometry-score differences were $-0.00331$, $-0.00451$
and $-0.00227$ at seeds 11, 23 and 37. The first seed motivated confirmation with
the remaining two. Their mean improvement was $0.00339$ (paired 95\% target bootstrap
$[0.00129,0.00590]$ on the same 128 validation targets). This modest change in a geometric
score has no measured full-test recovery counterpart. The primary benchmark keeps the
declared head.

\paragraph{Cost of the training representation.}
On one CPU thread, for a fixed FP32 batch of 32 training structures, the median of
23 warmed calls was 2.33~ms to enumerate indices and construct edge vectors, and
8.29~ms for a complete forward generator evaluation including its own neighbors.
Feature construction alone took 4.10~ms with pair shells and 7.81~ms with all powers.
The complete generator, descriptor and backward pass took 21.72 and 32.03~ms, respectively.
Summing the storage sizes seen by autograd's saved-tensor hook, with repeated storages
counted on each save, gave 6.05 and 9.11~MB, respectively. Process peak memory and
end-to-end proposal latency are measured separately in the benchmark.

\marksection{csp}{Prediction from composition}
\section{Predicting structures from composition}\label{sec:formula-csp}
The official MP-20 split contains 27,136 training, 9,047 validation and 9,046 test records. The earlier fixed-basis study had already examined the test set. Each model receives the elemental multiset in the unit cell. Eight uniformly scheduled training rows per update supply labels from 24,830 exact unit-cell compositions. For each occurrence we draw four observed structures independently with replacement, four independent generated negatives and four independent queries. A full scheduling pass visits every training row once. For 86.01\% of scheduled rows, the exact unit-cell composition has only one training structure. Sampling this bank repeatedly reproduces that structure. Across all three complete schedules, 70,922 of 81,408 positive banks contained one distinct training ID, 7,384 contained two, 2,274 contained three and 828 contained four.

A ridge regression estimates log volume per site from the elemental fractions and a site-count indicator. We add a Gaussian residual fitted for that count and truncate the result at the minimum and maximum observed training log volumes per site. The six traceless log-stretch entries follow a count-specific Gaussian. We clip its principal log stretches to $[-1.2,1.2]$ and subtract their mean. The canonical cell in the training basis and independent uniform sites form the physical seed. A zero-call control returns that seed.

Pair-shell, center-power-zero and full angular drift share the same physical backbone, random initial weights, training records and noise. Center power zero retains radial cross-products at individual centers. Orders one and two add dependence on the bond angle. We also train a fixed-basis set network with the same angular feature map. Its bounded residual updates the six log-stretch entries and fractional sites. A separate compact flow receives the same physical seeds and examples. Its velocity interpolates in the log-stretch chart and along mean-free wrapped site displacements after a random permutation within each species. It generates in sixteen Euler evaluations. In the center-power implementation we computed the full descriptor and sliced away orders one and two before evaluating the field. Physical-arm inference uses the same network operations at every angular order.

A deterministic 24-update fit to one training structure (\texttt{mp-1078776}) reduced squared descriptor error from 2.457 to 0.835, establishing that the query path could move the geometry. For the distributional fits we used AdamW at $10^{-3}$, clipped the gradient norm at ten, and fitted seeds 11, 23 and 37. After a validation pilot we compared 256, 1,024 and 3,392 updates, each with 32 observed examples. For each objective we selected the lowest mean validation energy score on four geometric summaries (log volume per atom and three nearest-distance statistics). All runs used five draws on the same 128 hash-selected validation IDs with the earlier training scales. The table gives these validation scores, summed process CPU hours over the three full trajectories, and the largest recorded process working set. Workers can run more than one fit, so the process high-water mark can include preceding work. The positive-feature cache required 8.91 s on one thread for 27,136 structures and occupied 259 MB as a float32 matrix. We kept the backbone at width 24 throughout the comparison.

We scored each ordered raw draw against its own target with the earlier pymatgen 2024.4.13 matcher. Length and site tolerances were 0.3 and 0.5, angle tolerance was $10^\circ$, and matching used primitive reduction, volume rescaling and exact species. RMS displacements are normalized by the matcher's characteristic per-site length. Neither ranking nor relaxation enters the proposal or matching cost. The paired intervals resample material IDs 5,000 times, taking their seeds and draws together. The earlier models retain the ordered test candidates from their released checkpoints. uFlowCSP contributes one checkpoint. Per-seed hit counts and the matched-target denominators are included in the tables accompanying the paper.

The GPU timing cohort comprises 128 hash-selected targets from the full multielement split. The earlier 128-target binary pilot is a separate cohort. The GPU comparisons used the same RTX 3060 Laptop GPU in FP32 with TF32 off, one Torch host thread and five warmups. Each serial candidate and each batch includes the prior, image lists, cell and site decoding, transfers and materialization of species, canonical lattice, fractional sites and Cartesian positions on the host. Host output synchronizes the GPU. Retrieval and the seed-only model run on the CPU. A separate batch-24 workload cycles the same compositions through exactly 10,000 requests. Loading, training-index preparation and retrospective scoring are outside the clock. Other CPU work shared this workstation.

\begin{center}\footnotesize\resizebox{\linewidth}{!}{\begin{tabular}{lrrrrrrr}\hline Objective & Parameters & 256 & 1,024 & 3,392 & Selected & CPU h & Peak RSS (MiB) \\ \hline
Physical, pair shells & 6,723 & 2.323 & 2.375 & 2.397 & 256 & 0.159 & 602 \\
Physical, center power 0 & 6,723 & 2.436 & 2.380 & 2.385 & 1,024 & 0.220 & 602 \\
Physical, angular 0--2 & 6,723 & 2.080 & 2.582 & 3.040 & 256 & 0.188 & 602 \\
Fixed basis, angular 0--2 & 6,497 & 2.691 & 2.758 & 2.768 & 256 & 0.178 & 602 \\
Compact flow (16, matched) & 6,497 & 2.853 & 2.741 & 2.703 & 3,392 & 0.024 & 602 \\
\hline\end{tabular}}\end{center}
\begin{figure}[!htbp]\centering\includegraphics[width=.94\linewidth]{csp-validation.pdf}\caption{Validation scores at the three matched update budgets. Thin lines show the three fits and heavy lines show their means. Open circles mark the frozen selections. All 45 points use the same 128 validation targets and five draws per target. Each full physical angular fit deteriorated by the final budget. Backbone width and optimizer family were fixed across the new fits.}\label{fig:v2-csp-validation}\end{figure}
\begin{table}[!htbp]\centering\small\caption{Raw CSP on all 9,046 MP-20 test targets. MR uses all targets. RMS averages the best accepted normalized displacement per recovered target within each seed, with equal seed weights. $S$ denotes neural evaluations per candidate. At 1, 5 and 20 draws each target uses $S,5S,20S$ evaluations. B1 is median warm batch-one latency.}\label{tab:v2csp}
\resizebox{\linewidth}{!}{\begin{tabular}{lrrrrrrrrrr}\hline & & & \multicolumn{3}{c}{MR (\%)}&\multicolumn{3}{c}{Mean matched RMS}&\multicolumn{2}{c}{Wall clock}\\
Model & $S$ & $kS$ & $k=1$ & 5 & 20 & 1 & 5 & 20 & B1 (ms) & min/10k\\ \hline
Fitted prior alone & 0 & 0/0/0 & 5.05 & 15.96 & 26.73 & 0.395 & 0.391 & 0.371 & 0.36 & 0.004\\
Physical, pair shells & 1 & 1/5/20 & 6.60 & 19.36 & 31.30 & 0.382 & 0.370 & 0.346 & 17.19 & 0.160\\
Physical, center power 0 & 1 & 1/5/20 & 6.70 & 19.36 & 31.47 & 0.382 & 0.369 & 0.346 & 15.06 & 0.110\\
Physical, angular 0--2 & 1 & 1/5/20 & 6.83 & 19.68 & 31.79 & 0.381 & 0.367 & 0.344 & 23.06 & 0.184\\
Fixed basis, angular 0--2 & 1 & 1/5/20 & 5.28 & 17.51 & 30.27 & 0.390 & 0.383 & 0.362 & 11.51 & 0.189\\
Compact flow (16, matched) & 16 & 16/80/320 & 7.35 & 21.12 & 33.56 & 0.373 & 0.358 & 0.334 & 97.03 & 0.679\\
Released fixed-basis drift & 1 & 1/5/20 & 6.79 & 20.33 & 32.91 & 0.382 & 0.371 & 0.346 & 2.83 & 0.022\\
Released compact flow (16) & 16 & 16/80/320 & 10.63 & 23.71 & 36.54 & 0.363 & 0.334 & 0.319 & 37.52 & 0.226\\
Compact MeanFlow & 1 & 1/5/20 & 8.76 & 21.41 & 31.86 & 0.388 & 0.362 & 0.339 & 3.18 & 0.025\\
Released uFlowCSP & 1 & 1/5/20 & 46.32 & 68.05 & 79.32 & 0.089 & 0.088 & 0.085 & 14.49 & 0.621\\
Train-only retrieval & 0 & 0/0/0 & 61.86 & 64.60 & 65.19 & 0.040 & 0.041 & 0.042 & 0.08 & 0.003\\
\hline\end{tabular}}\end{table}
\begin{figure}[!htbp]\centering\includegraphics[width=\linewidth]{csp-recovery.pdf}\caption{Recovery on the same 128 hash-selected MP-20 targets at equal completed counts and equal elapsed proposal time. Each target owns a serial stream capped at 20 proposals. The shaded pointwise 95\% interval resamples IDs for the full angular generator. Source tables include intervals for every curve. The vertical line marks 100 ms.}\label{fig:v2csp}\end{figure}
On the validation energy score, angular minus pair-shell at equal exposure was -0.243 \([-0.264,-0.221]\) after 256 updates and +0.643 \([+0.538,+0.750]\) after 3,392. Lower scores favor the first method. Thus the early angular advantage reversed by the full scheduling pass.

One Euler update of the earlier compact velocity checkpoint is a separate diagnostic. With its original same-target candidates, MR at 1/5/20 draws was 7.49/18.00/25.04\% and mean matched normalized RMS was 0.403/0.389/0.359. Its latency belongs to the earlier timing stream.

At twenty draws, angular training changed the recovered fraction by \(+0.49\;[+0.38,+0.60]\) percentage points relative to pair-shell drift and by \(+0.32\;[+0.15,+0.49]\) relative to center power alone. Brackets give paired 95\% intervals. The change from the fitted prior was \(+5.06\;[+4.74,+5.37]\). With the same angular field, changing from the fixed-basis backbone to the physical network gave \(+1.51\;[+1.11,+1.92]\) points. Relative to the earlier fixed-basis release, the change was \(-1.12\;[-1.54,-0.71]\).

Angular and pair-shell drift both selected 256 updates. Their test contrast isolates the feature change at the same update count. The center-power control selected 1,024 updates, so its test contrast also includes the selected training duration.

At twenty draws the angular matched-RMS denominators in seeds 11, 23 and 37 were 2,860, 2,873, 2,893. Its mean matched RMS was 0.3436 \([0.3402,0.3470]\). Because recovery changes with the method and draw count, these conditional displacement averages can describe different target populations.

Two stricter checks on the same angular draws give 17.81\% recovery with the maximum-site matcher and 21.24\% when the accepted RMS draw must also have a native-volume ratio in $[1/1.10,1.10]$, both at twenty draws.

In the sequential benchmark at 100 ms, the angular model recovered 18/128 targets with 4.09 completed draws per target. The fitted prior recovered 36, pair-shell drift 20, released fixed-basis drift 39, matched iterative flow 5 and uFlowCSP 90.

The seed-only validation energy score averaged 3.159. The fifteen full training trajectories consumed 0.77 CPU worker-hours and zero GPU hours. The archive records 3,256,560 new candidates, including 0 matcher exceptions and 522 direct public matcher checks.

Peak GPU allocations in the complete proposal benchmark were 11.6 MiB for angular drift, 11.6 MiB for pair-shell drift and 601.5 MiB for uFlowCSP. These include the resident model and timed activations. The training descriptor is absent from each physical generator's sampling path.

On one Intel Core i7-12700H thread, median warm latencies for the fitted prior, pair-shell drift, angular drift and matched sixteen-step flow were 0.45, 4.55, 4.57, 17.14 ms. Their separate batch-24 workloads of exactly 10,000 host-decoded outputs took 0.28, 4.73, 5.29, 12.34 s in the same order. The 128 CPU targets each received five batch-one calls per method, with method order balanced within six-target blocks. Concurrent matching jobs shared the CPU during these clocks.

We repeated the serial deadlines after balancing method order within six-target blocks (Figure~\ref{fig:v2balanced}). All six services remained resident. Each target--method pair ran for twenty consecutive draws, with shared per-target noise for the physical arms. At 100 ms, the fitted prior, pair-shell drift and angular drift recovered 39, 36 and 32 of 128 targets. The earlier fixed-basis generator recovered 38, the matched sixteen-step flow 13, and the earlier sixteen-step flow 13. The evidence package retains both streams, their separate proposals and paired target intervals.


\paragraph{Published operating points.}
Dey et al.\ report 78.38\% recovery and mean normalized RMS 0.0783 for one-step uFlowCSP with 20 proposals per MP-20 target. Five steps at the same draw count reached 83.64\% and 0.0619. Their 10,000-proposal clocks on an A100 were 0.39 and 1.31 minutes \cite{uflow}. Their comparison lists 77.93\% for 1,000-step DiffCSP and 78.34\% for 100-step CrystalFlow at 20 draws \cite{diffcsp,crystalflow}. In the fixed-basis release we separately ran the official CrystalFlow composition checkpoint at five unannealed Euler steps and recovered 7,768 of the 9,046 targets at 20 draws (85.87\%). Those saved proposals belong to the released CrystalFlow checkpoint. All timing comparisons for the new physical networks use our local measurements.
\marksection{dng}{Generation across multielement compositions}
\section{De novo generation with matched sampled chemistry}\label{sec:denovo}
We tested the geometry learned from formula-conditioned prediction in a complete sampling pipeline. A train-only prior draws a cell multiplicity pattern and then its distinct atomic numbers in ascending order. Within each pattern the next-element counts shrink toward a shared element-rank distribution; the shrinkage strength of four was selected on validation in the earlier release. The prior receives no target formula. Its random stream produced the same 10,000 compositions in all three arms, including 6,478 draws at reduced formulas absent from MP-20 training. The physical angular and radial networks share the physical seed prior, fit seed and inference noise.
\begin{table}[!htbp]\centering\small\setlength{\tabcolsep}{3.8pt}
\caption{De novo geometry on 10,000 unfiltered draws per arm. Structural (SV), compositional (CV) and joint validity (JV) use all attempts. Coverage recall uses the 9,046 test structures; precision uses all 10,000 generated structures. Density, formation energy and distinct-element Wasserstein distances use 1,000 valid draws selected independently within each arm.}\label{tab:new-raw}
\begin{tabular}{lrrrrrrrr}\hline Method & SV (\%) & CV (\%) & JV (\%) & COV-R (\%) & COV-P (\%) & $W_\rho$ & $W_E$ & $W_{\rm el}$\\\hline
Fixed-basis release & 90.70 & 83.57 & 75.84 & 49.83 & 99.09 & 0.747 & 0.378 & 0.049\\
Physical radial & 94.29 & 83.57 & 78.83 & 60.84 & 99.04 & 0.577 & 0.354 & 0.068\\
Physical angular & 95.59 & 83.57 & 79.88 & 60.48 & 99.03 & 0.276 & 0.373 & 0.057\\
\hline\end{tabular}\end{table}
Following the released DiffCSP evaluator, we require the minimum distinct-site separation to be at least 0.5~\AA\ and the cell volume to be at least 0.1~\AA$^3$ \cite{diffcsp}. We test charge/electronegativity compatibility with SMACT, with the released elemental and all-metal exceptions \cite{smact}. Joint validity also requires successful fingerprints. The structure fingerprint is the site-averaged 61-component CrystalNN vector. The 132-component Magpie composition vector is standardized with the released constants \cite{matminer}. Coverage requires a neighbor within both thresholds, 0.4 in structure space and 10 in chemistry space; the minimizing neighbors may differ. Density is in g\,cm$^{-3}$. The released MP-20 property network supplies the raw formation-energy values in eV/atom. Each distance uses the same predictor or observable on all 9,046 references.
\begin{table}[!htbp]\centering\small\setlength{\tabcolsep}{4pt}
\caption{Raw diversity and measured end-to-end CPU cost. U is the first representative in its generated bank, N has no training-set structural match, and V denotes joint validity. Each count is out of 10,000 attempts. Time includes composition and geometry draws, graph construction, and decoding to host coordinates. All three generators use one network evaluation per attempt.}\label{tab:new-cost}
\begin{tabular}{lrrrrrr}\hline Method & U & N & V$\cap$U$\cap$N & B1 (ms) & 10,000 (s) & V$\cap$U/s\\\hline
Fixed-basis release & 9,943 & 9,856 & 7,408 & 1.44 & 2.27 & 3322.6\\
Physical radial & 9,944 & 9,857 & 7,703 & 6.46 & 8.27 & 947.1\\
Physical angular & 9,942 & 9,850 & 7,800 & 5.96 & 8.29 & 956.8\\
\hline\end{tabular}\end{table}
The angular-minus-radial difference in structural validity was +1.30 percentage points (paired 95\% interval $[1.08,1.53]$). Coverage recall changed by -0.36 points ($[-0.83,0.08]$). The first interval resamples generated draw pairs; the second resamples test-reference coverage indicators conditional on the two realized banks. Each result uses a single fitted seed. The raw novelty check applies a species-aware primitive-cell match against all 27,136 training structures with the same reduced formula as the draw.
\paragraph{Relaxation and a restricted competing-phase hull.} The first sixteen draws of each new arm entered the same CHGNet 0.3.0 zero-pressure cell relaxation using ASE FIRE for at most 100 steps and a 0.10~eV/\AA\ generalized force tolerance \cite{chgnet}. For each chemical system, the MP-20 near-hull library includes the lowest deposited formation-energy example per available reduced formula at or below 0.10~eV/atom on the deposited hull, plus elemental references. We use the same relaxed CHGNet energy convention for candidates and competing phases. Every phase and elemental terminal must be available for a complete label. The candidate-augmented hull takes the nonnegative part of the energy difference to the best competing mixture. A screened hit must converge, avoid contacts below 0.60 times the sum of covalent radii, lie within 0.10~eV/atom of that restricted hull, be unique within its relaxed cohort and have no MP-20 training match.
The radial model had 8/16 converged relaxations, 4 final short-contact cases, 16/16 complete hulls and 0/16 screened hits. Candidate relaxations took 258.7 CPU worker-seconds, with 0.12 seconds for final identity and hull checks.
The angular model had 6/16 converged relaxations, 4 final short-contact cases, 16/16 complete hulls and 0/16 screened hits. Candidate relaxations took 232.0 CPU worker-seconds, with 0.24 seconds for final identity and hull checks.
The 195 distinct reference phases shared by the two arms all converged. Their charged cost was 238.5 CPU worker-seconds per arm. Charging generation, candidate and phase relaxations, and final scoring gives zero screened hits per summed worker-hour in each prospectively selected sixteen-draw cohort. A missing competing phase could alter a restricted-hull energy. The present physical-periodic checkpoints have no DFT-confirmed discovery yield. The eight earlier single-point calculations belong to the fixed-basis study.
\begin{figure}[!htbp]\centering\includegraphics[width=.78\linewidth]{new-crystals.pdf}
\caption{Generated and unconstrained CHGNet-relaxed cells on the same physical scale. We traverse the angular draws and then the radial draws in their original orders, selecting the first cell with at least three species to converge with acceptable contacts and the first to fail either gate, wherever available. Intracell sites are drawn in the deposited basis. The accompanying CIFs contain the complete coordinates.}\label{fig:denovo-examples}\end{figure}

\paragraph{Which geometries did the prior reach?} The shared chemistry sample contains 3,522 draws at seen reduced formulas and 6,478 at unseen ones. The first set assesses geometry where training contains at least one member of that formula. The second tests extrapolation to previously unseen stoichiometries.
The radial generator passed structural validity on 3,323/3,522 seen draws and 6,106/6,478 unseen draws. In the first sixteen proposals its relaxations converged on 3/4 and 5/12, respectively.
The angular generator passed structural validity on 3,365/3,522 seen draws and 6,194/6,478 unseen draws. In the first sixteen proposals its relaxations converged on 3/4 and 3/12, respectively.
\paragraph{Earlier experiments with the same evaluator.} Using the empirical training-composition prior, the earlier fixed-basis checkpoint generated no unseen formulas and attained 61.67\% coverage recall. Using one-step uFlowCSP on the identical learned-prior chemistry, our earlier replay attained 99.47\% recall. Our separate replay of joint CrystalFlow sampled its own chemistry and geometry with 100 Euler steps, generating 8,550 training-unseen formula draws among 10,000, 82.19\% joint validity and 98.80\% coverage recall. That bank took 29.33 minutes on the same GPU under its own batch-24 protocol. Six of its first sixteen proposals passed the restricted physical screen. Formula overlap, raw geometry, and screened physical yield measure different stages of the generator \cite{uflow,crystalflow}.

\paragraph{Common GPU workload.} On the RTX 3060 Laptop GPU, fixed-basis, radial and angular generation took 1.25, 4.92, 3.10 seconds for 10,000 structures at batch 128. Their warm batch-one medians were 3.93, 17.82, 19.29 ms. The clock includes drawing the composition, sampling a physical seed, decoding one network evaluation, transferring to host memory and constructing Cartesian coordinates. It ends before fingerprinting or relaxation. Applying the corresponding CPU-scored valid-unique counts to these clocks gives estimated yields of 361,026, 95,504, 153,427 per GPU minute. CPU and GPU streams had identical chemistry, with worst case differences of $6.5\times10^{-7}$ in a log-metric coordinate and $3.6\times10^{-7}$ in a wrapped fractional coordinate across the three paired 10,000-cell replays.

\marksection{symmetry}{Prediction from occupied Wyckoff sites}
\section{Generation with a supplied occupied Wyckoff template}\label{sec:wyckoff}
The second task supplies a Hall setting and the multiset of element-labelled occupied Wyckoff letters. A letter can occur more than once. These inputs determine conventional-cell site counts, but provide neither the cell volume nor any free coordinate. We use the first 128 official test identifiers in their SHA-256 ordering. Their expanded conventional cells contain up to 60 sites. For each occupied orbit, we independently sample its free coordinates, expand the tabulated coset images and draw a compatible metric from a train-only prior. We evaluate a frozen formula-trained angular map on these seeds and fit a second angular map on expanded conventional crystals conditional on their supplied template. These are protocol C results. Protocol A supplies only the unit-cell composition.

\begin{table}[!htbp]\centering\footnotesize\setlength{\tabcolsep}{3.4pt}
\caption{Prediction from the same supplied Wyckoff templates on 128 test targets. MR (\%) and mean best normalized RMS use ordered raw draws. RMS averages the matched targets. Valid (\%) is the structural-validity fraction among all draws. Retained is the number of the first 16 targets whose first draw retains the requested subgroup after unconstrained CHGNet relaxation at 0.01\,\AA. The native DiffCSP++ row evaluates one draw.}\label{tab:wyckoff}
\resizebox{\linewidth}{!}{\begin{tabular}{lrrrrrrrrr}\hline
 & & \multicolumn{3}{c}{MR} & \multicolumn{3}{c}{RMS}& & \\
Method & NFE & 1 & 5 & 20 & 1 & 5 & 20 & Valid & Retained\\\hline
Template prior & 0 & 50.8 & 70.3 & 78.9 & 0.180 & 0.187 & 0.160 & 82.2 & 12\\
A angular, template seeds & 1 & 56.2 & 74.2 & 82.8 & 0.168 & 0.165 & 0.139 & 89.2 & 14\\
C angular & 1 & 56.2 & 73.4 & 82.8 & 0.171 & 0.164 & 0.141 & 88.9 & 15\\
C angular + orbit variance & 1 & 56.2 & 74.2 & 82.8 & 0.168 & 0.164 & 0.139 & 89.6 & 14\\
DiffCSP++ (respaced) & 40 & 82.0 & 92.2 & 94.5 & 0.044 & 0.032 & 0.026 & 100.0 & 16\\
DiffCSP++ (native) & 1,998 & 86.7 & -- & -- & 0.031 & -- & -- & 100.0 & 16\\
\hline\end{tabular}}\end{table}
The C-trained angular map recovered 106 targets at twenty draws, compared with 101 for its train-fitted metric prior and random Wyckoff coordinates used alone. The paired recovery difference was 3.91 percentage points (95\% target-bootstrap interval $[0.78,7.03]$). Adding the variance across repeated orbits recovered the same 106 targets. Their mean matched RMS values were 0.141 and 0.139.
DiffCSP++ supplies an independent predictor using the same group, letters and species \cite{diffcspplus}. We replayed its released MP-20 weights with 20 predictor--corrector iterations (40 neural calls) at all three draw budgets and repeated the first-draw comparison with the complete 999-iteration sampler (1,998 calls). The abbreviated sampler preserves the authors' coordinate projections and indexes their learned noise predictions at 20 equally spaced original time indices. It uses the corresponding noise marginals and respaced lattice posterior variance. SGFM also constrains each trajectory by supplied orbit information \cite{sgfm}.

Across all 9,046 official test records, one draw from the template prior alone recovered 4,152 targets (45.90\%), and one C-angular proposal recovered 4,549 (50.29\%). The seven rejected template round trips remain in the denominator. The paired angular-minus-prior difference was 4.39 points ($[3.90,4.90]$).
At 0.01\,\AA, the explicit site-orbit and metric tests accepted all 2,560 C-angular draws. Independent symmetry finding contained the requested operations on 2,557 draws and assigned the exact requested group type on 2,556. The detected operation set equalled the requested set in its supplied setting on 2,556 draws. After relaxation, 15/16 retained the requested subgroup, 15/16 recovered its operation set and 15/16 returned the requested group type. The raw guarantee concerns the operations shared by a template-compatible seed. Additional symmetry can occur, and unconstrained relaxation can alter the group.
\begin{figure}[!htbp]\centering\includegraphics[width=.94\linewidth]{wyckoff-retention.pdf}
\caption{Requested subgroup containment, exact finite operation set in the supplied setting, and exact international group type at 0.01\,\AA. Error bars give 95\% bootstrap intervals resampling whole target IDs. The left panel expands the range from 99.6 to 100\% and includes twenty draws on each of 128 targets. The right panel follows the first draw from the first sixteen targets per method. An optimizer step cap retains its last observable geometry.}\label{fig:wyckoff-retention}
\end{figure}

\section{Discussion}
A representation-independent training field solves a concrete problem. A cell can be redrawn in a skew basis while another member of the same bank keeps its original basis. The descriptor then gives the same displacement and the same parameter gradient. This statement survives restoration of the trained networks. Its limits are visible in small examples. Preserving which radial shells meet at a common atom already separates the two four-site motifs. Truncating the neighborhood still erases longer-range order, and even-parity contractions equate mirror images.

Standard twenty-draw recovery increased by 0.49 percentage points when angular features replaced pair shells on the same backbone. Native volume changed that conclusion. Requiring recovery and agreement with the reference cell volume reduced angular recovery to 21.24\%, versus 24.60\% for pair shells (Section~\ref{app:csp-sensitivity}). The scale-enabled matcher was rewarding site arrangements that the angular model placed in cells of less accurate native volume.

At 100 ms in the balanced-order replay, the angular map recovered 32 targets against 13 for the sixteen-call flow, a gain of 14.84 percentage points (paired 95\% interval 8.59--21.09). Radial training recovered 36, the earlier fixed-basis generator recovered 38, and the fitted seed alone recovered 39. One graph update takes enough time to change which model leads at a fixed deadline. The table of full-split scores measures a different use of the sampler, in which each target receives its full allowance of completed draws.

What the generator learns depends strongly on the seed. Most exact cell compositions in these training schedules have one observed structure. Sampling a larger positive bank for one of these labels repeats that observation. The compact head can only make a bounded deformation in one call, and longer angular training worsened the validation geometry score. On an occupied Wyckoff template, the seed already fixes multiplicity, permitted coordinates and lattice symmetry. The learned map improves recovery over this seed on the full split. On the matched 128-template cohort, adding orbit variance left twenty-draw recovery unchanged, while the released diffusion sampler placed atoms more precisely.

The de novo results distinguish geometric feasibility from physical quality. On identical chemistry and noise, angular training raised structural validity from 94.29\% to 95.59\%, yet its coverage recall was 60.48\% against 60.84\% for radial training. The validity gain thus coexists with similar coverage of the held-out geometries. Force convergence and competition with phases of the same chemical system imposed stronger tests on the first sixteen proposals. That cohort is small and the reference hull uses a restricted set of MP-20 phases and a learned potential. A discovery yield for these checkpoints would require broader competing phases and independent electronic-structure calculations. The evidence supports one-call generation with a proven symmetry and a measured computational trade-off. Its strongest practical use in these experiments is refining a supplied symmetry-compatible seed.

\paragraph{Data, code and authorship.}
The accompanying research package contains the checkpoint and split identifiers, ordered draws, evaluator settings, editable figures and run commands. The earlier 38-page release remains the fixed-basis benchmark. The official MP-20 test structures were inspected in that work. The present model choices use training and validation. OpenAI Codex assisted with literature verification, software, experiments, analysis and drafting.
\begin{thebibliography}{99}\small\setlength{\itemsep}{1.5pt}
\bibitem{deng} M. Deng, H. Li, T. Li, Y. Du and K. He. \linkto{https://arxiv.org/abs/2602.04770}{Generative Modeling via Drifting}. arXiv:2602.04770 (2026).
\bibitem{symdrift} S. Darouich et al. \linkto{https://arxiv.org/abs/2605.06140}{SymDrift: One-Shot Generative Modeling under Symmetries}. arXiv:2605.06140 (2026).
\bibitem{soap} Albert P. Bartók, Risi Kondor and Gábor Csányi. \linkto{https://journals.aps.org/prb/abstract/10.1103/PhysRevB.87.184115}{On representing chemical environments}. Phys. Rev. B 87, 184115 (2013).
\bibitem{whole} Sandip De, Albert P. Bartók, Gábor Csányi and Michele Ceriotti. \linkto{https://pubs.rsc.org/en/content/articlelanding/2016/cp/c6cp00415f/unauth}{Comparing molecules and solids across structural and alchemical space}. Phys. Chem. Chem. Phys. 18, 13754--13769 (2016).
\bibitem{degenerate} Sergey N. Pozdnyakov et al. \linkto{https://journals.aps.org/prl/abstract/10.1103/PhysRevLett.125.166001}{Incompleteness of Atomic Structure Representations}. Phys. Rev. Lett. 125, 166001 (2020).
\bibitem{lego} Osman Goni Ridwan et al. \linkto{https://www.nature.com/articles/s41524-025-01931-9}{AI-assisted rapid crystal structure generation towards a target local environment}. npj Comput. Mater. 12, 64 (2026).
\bibitem{kernelgrad} M. Esteban-Casadevall et al. \linkto{https://arxiv.org/abs/2605.10727}{Kernel-Gradient Drifting Models}. arXiv:2605.10727 (2026).
\bibitem{tokendrift} D. Oba, H. Furuta and N. Okazaki. \linkto{https://arxiv.org/abs/2605.19470}{Drifting Objectives for Refining Discrete Diffusion Language Models}. arXiv:2605.19470 (2026).
\bibitem{driftreact} R. Schlama and P. Schwaller. \linkto{https://arxiv.org/abs/2605.22990v2}{Drift-React: One-step Generation of Reaction Pathways via SE(3) Drifting Fields}. arXiv:2605.22990v2 (2026).
\bibitem{egnn} Víctor Garcia Satorras, Emiel Hoogeboom and Max Welling. \linkto{https://proceedings.mlr.press/v139/satorras21a.html}{E(n) Equivariant Graph Neural Networks}. Proc. ICML, PMLR 139, 9323--9332 (2021).
\bibitem{gems} Astrid Klipfel, Yaël Frégier, Adlane Sayede and Zied Bouraoui. \linkto{https://www.ijcai.org/proceedings/2023/669}{Unified Model for Crystalline Material Generation}. Proc. IJCAI, 6031--6039 (2023).
\bibitem{vector} Astrid Klipfel, Yaël Fregier, Adlane Sayede and Zied Bouraoui. \linkto{https://ojs.aaai.org/index.php/AAAI/article/view/30224}{Vector Field Oriented Diffusion Model for Crystal Material Generation}. Proc. AAAI 38(20), 22193--22201 (2024).
\bibitem{crystalflow} X. Luo et al. \linkto{https://doi.org/10.1038/s41467-025-64364-4}{CrystalFlow: a flow-based generative model for crystalline materials}. Nature Communications (2025).
\bibitem{flowmm} B. K. Miller, R. T. Q. Chen, A. Sriram and B. M. Wood. \linkto{https://proceedings.mlr.press/v235/miller24a.html}{FlowMM: Generating Materials with Riemannian Flow Matching}. ICML (2024).
\bibitem{diffcsp} R. Jiao et al. \linkto{https://proceedings.neurips.cc/paper_files/paper/2023/hash/38b787fc530d0b31825827e2cc306656-Abstract-Conference.html}{Crystal Structure Prediction by Joint Equivariant Diffusion}. NeurIPS (2023).
\bibitem{uflow} S. Dey et al. \linkto{https://arxiv.org/abs/2609.09799}{uFlowCSP: Crystal Structure Prediction using Mean flow generative models}. arXiv:2609.09799 (2026). \linkto{https://github.com/usccolumbia/uFlowCSP}{Implementation and checkpoint links}.
\bibitem{omatg} T. Egg, H. W. Sullivan, E. B. Tadmor and S. Martiniani. \linkto{https://arxiv.org/abs/2609.26402}{OMatG-flash: An All-Atom Flow Map with Reinforce Adjoint Matching for Scalable Materials Discovery}. arXiv:2609.26402 (2026).
\bibitem{smact} D. W. Davies et al. \linkto{https://doi.org/10.21105/joss.01361}{SMACT: Semiconducting Materials by Analogy and Chemical Theory}. Journal of Open Source Software 4, 1361 (2019).
\bibitem{matminer} L. Ward et al. \linkto{https://doi.org/10.1016/j.commatsci.2018.05.018}{Matminer: An open source toolkit for materials data mining}. Computational Materials Science 152, 60--69 (2018).
\bibitem{chgnet} B. Deng et al. \linkto{https://doi.org/10.1038/s42256-023-00716-3}{CHGNet as a pretrained universal neural network potential for charge-informed atomistic modelling}. Nature Machine Intelligence 5, 1031--1041 (2023).
\bibitem{diffcspplus} R. Jiao et al. \linkto{https://arxiv.org/abs/2402.03992}{Space Group Constrained Crystal Generation}. arXiv:2402.03992 (2024).
\bibitem{sgfm} O. Puny et al. \linkto{https://arxiv.org/abs/2509.23822}{Space Group Conditional Flow Matching}. arXiv:2509.23822 (2025).
\bibitem{wyformer} N. Kazeev et al. \linkto{https://proceedings.mlr.press/v267/kazeev25a.html}{Wyckoff Transformer: Generation of Symmetric Crystals}. ICML (2025).
\bibitem{matra} P. De Breuck et al. \linkto{https://arxiv.org/abs/2501.16051}{A generative material transformer using Wyckoff representation}. arXiv:2501.16051 (2025).
\bibitem{st2m} J. Liu et al. \linkto{https://arxiv.org/abs/2611.02906}{Site-orbit masked twin-branch Transformer for symmetry-constrained crystal generation}. arXiv:2611.02906 (2026).
\bibitem{dynacrys} Z. Jin et al. \linkto{https://arxiv.org/abs/2608.07401}{DynaCrys: Crystal Generation with Dynamic Space-Group Diffusion}. arXiv:2608.07401 (2026).
\bibitem{sgequi} R. Chang et al. \linkto{https://arxiv.org/abs/2505.10994}{Space Group Equivariant Crystal Diffusion}. arXiv:2505.10994 (2025).
\bibitem{orbitol} I. \v{Z}ugec, J. Rosenthal, I. Lon\v{c}ari\'c, J.-W. van de Meent and T. Had\v{z}i Veljkovi\'c. \linkto{https://arxiv.org/html/2610.30086v1}{Orbitol: Crystal Generation on Wyckoff Orbits}. arXiv:2610.30086 (2026).
\bibitem{mattergen} C. Zeni et al. \linkto{https://doi.org/10.1038/s41586-025-08628-5}{A generative model for inorganic materials design}. Nature 639, 624--632 (2025). \linkto{https://github.com/microsoft/mattergen/tree/92423660a8bd70e83679086e88f88596d484dc16}{Alex-MP-20 source archive}.
\bibitem{pyxtal} Fredericks et al. \linkto{https://doi.org/10.1016/j.cpc.2020.107810}{PyXtal: A Python library for crystal structure generation and symmetry analysis}. (2021).
\bibitem{spglib} Togo et al. \linkto{https://doi.org/10.1080/27660400.2024.2384822}{Spglib: a software library for crystal symmetry search}. (2024).
\bibitem{zrelation} J. Mandereau, D. Ghisi, E. Amiot, M. Andreatta and C. Agon. \linkto{https://doi.org/10.1080/17459737.2011.608819}{Z-relation and homometry in musical distributions}. Journal of Mathematics and Music 5 (2011).
\bibitem{homometry} U. Grimm and M. Baake. \linkto{https://arxiv.org/abs/0808.0094}{Homometric Point Sets and Inverse Problems}. Zeitschrift f\"ur Kristallographie 223, 777--781 (2008).
\bibitem{lang2str} C. Liu, C. Gong, Z. Liu, J. Zhao and Y. Zhang. \linkto{https://arxiv.org/abs/2603.03946}{Lang2Str: Two-Stage Crystal Structure Generation with LLMs and Continuous Flow Models}. arXiv:2603.03946 (2026).
\end{thebibliography}

\clearpage
\appendix
\section*{Mathematical and experimental supplement}
\addcontentsline{toc}{section}{Mathematical and experimental supplement}
\marksection{supp}{Mathematical and experimental supplement}
\section{Coordinates, finite banks and maps of measures}
\paragraph{Wrapping and image indices.}
Write $F'=PFU^{-1}+\mathbf 1t^\top-W$ with integer wrap rows $W_i$. If the old edge was $i\to j$ with image $n$ and the new labels are $i',j'$, then its new integer index is
\[
 n'=nU^{-1}+W_{j'}-W_{i'} .
\]
Indeed $f'_{j'}-f'_{i'}+n'=(f_j-f_i+n)U^{-1}$ and right multiplication by $ULQ$ gives $dQ$. The inverse index map is integer-valued because $U$ and $U^{-1}$ are integer matrices. It covers coincident distinct atoms and preserves the condition of being a nonzero self-image. A primitive-to-conventional change that increases the cell multiplicity generally lacks this bijection on the indexed finite motif.

\paragraph{What a parameter gradient represents.}
For row perturbations at fixed lattice, $\delta x=\delta f\,L$. The pairing with a physical gradient $g_x$ yields $g_f=g_xL^\top$. A physical steepest-descent step written in fractional coordinates is $-g_f(LL^\top)^{-1}$, while the actual decoder learns a Cartesian displacement and changes its units using the seed lattice inverse. At parameters where the feature map is differentiable, two equivalent seed descriptions give the same scalar stopped-target loss for \emph{all} nearby parameter choices with the target fixed. The chain rule then equates the parameter gradients. At a graph cutoff the radial windows vanish with their first derivatives. Fractional wraps are chart cuts; a periodic physical observable has the same limiting value at both sides.

\paragraph{Independent finite banks.}
Condition on all query features, the model, and the list of chemical labels. For each block $c$ let $P_c$ and $Q_c$ be its positive and negative feature laws. The coordinates of the positive bank are independent samples from $P_c$; negative samples are independent of them and of the query noise. If $P_c=Q_c$ in every block, simultaneous exchange of the two banks preserves their joint law. Because the denominator in Eq.~\eqref{eq:v2field} averages across blocks, the simultaneous exchange is sufficient even when a calculation is implemented as a single batched reduction. A bank with a designated positive geometry fixed by its identity need not satisfy this conditional exchange. A singleton empirical formula gives the same positive point in every slot.

For $M$ query rows, each single-temperature normalized vector has norm at most $\sqrt M$, since its squared norm contributes one of the $M$ terms in the RMS. All finite sums, affinities and row-column normalizers are continuous, and the RMS denominator is strictly positive. If the finitely many conditional negative feature laws converge weakly to their corresponding positive laws, the finite product of their laws also converges weakly. Applying bounded continuous integration componentwise gives convergence of the conditional mean field to zero. A single realized drift need not tend to zero at fixed bank size. These statements concern laws after applying the chosen descriptor.

\paragraph{Finite structural support.}
Fix a nonsingular seed lattice, a formula with at least two sites and finite network parameters. In the ideal real-arithmetic map the fractional output lifts to $F\mapsto F+s(F)$ on $\mathbb R^{3n}$. The displacement $s(F)=v(F)L^{-1}$ is bounded, continuous and periodic in each input site. Edges that enter or leave the finite graph have vanishing weights at the cutoff. Choose $b$ bounding all components. For any specified output $y$, the continuous map $F\mapsto y-s(F)$ maps the closed cube $y+[-b,b]^{3n}$ into itself. Brouwer's theorem gives a fixed point, hence a preimage of $y$. On the torus the inverse image of every nonempty open set is thus a nonempty open set, which has positive probability under independent uniform input sites. Projection to the relative motif removes the common origin and still assigns positive mass to each open relative-coordinate set. This establishes full support for the fractional positional marginal at that seed lattice. It prevents equality to a finite list of native relative motifs, even though mass can concentrate near a few motifs. The resulting \emph{feature} marginal can lose information. The argument supplies no lower bound on the probability of a useful structure. Template seeds instead have support determined by their occupied charts.

\paragraph{Species channels and physical information.}
An orthogonal frame change acts on each order-$\ell$ coefficient vector by the corresponding real harmonic representation. Under inversion its factor is $(-1)^\ell$; two factors of the same order appear in the power, so mirror motifs remain indistinguishable. If handedness is physically relevant, two enantiomorphs inequivalent under proper rotations may thus have identical features. For a single center and neighbor channels $u,v$, the addition theorem expands a power component into a sum over ordered pairs of edges. The angular factor is proportional to $P_\ell(\widehat d_e\cdot\widehat d_{e'})$. Already at $\ell=0$ the two radial weights share a center. Pooling before multiplication would include terms from different centers. This is why the order-zero control is necessary to interpret the 12-gon experiment.

\paragraph{Conditional and unconditional draws.}
Let $\xi$ be a physical seed drawn in the declared chart for a supplied cell composition $c$. The formula-only law is the pushforward $Q_{\theta,c}=(G_\theta)_\#\mu_c$. The matched de novo law is $\pi(\mathrm dc)Q_{\theta,c}(\mathrm dx)$, where $\pi$ samples the composition using training counts. The endpoint receives neither a reference formula from the test set nor a test structure. For a supplied group and species-labelled template $T$, the analogous pushforward uses $\mu_{c,T}$. Its support lies in the subset fixed by every required operation, subject to the group construction and numerical tolerances. The number of operations fixing an individual output may be larger. A measured conditional success rate on supplied $T$ does not describe how often a model would propose $T$ on its own. A group-only input with an inferred template (protocol B) would require choosing orbit occupations. A generator that first sampled composition, group and occupied template (protocol D) would also require a joint prior for these labels. Our de novo study samples composition and predicts unrestricted geometry. Its space-group distribution is implicit.

\section{Scope of literature comparisons}
WyFormer models Wyckoff occupations; Matra-Genoa encodes discrete and continuous quantities in an autoregressive sequence; ST2M masks the free coordinates of occupied sites; and DynaCrys allows the symbolic symmetry to evolve \cite{wyformer,matra,st2m,dynacrys}. DiffCSP++, SGEquiDiff, SGFM and Orbitol provide closer precedents for generation at specified crystallographic symmetry \cite{diffcspplus,sgequi,sgfm,orbitol}. Their input settings determine whether any recovery comparison is meaningful. In particular, inherited occupied sites and labelled elements considerably constrain a target. MatterGen jointly generates chemistry and geometry, with its own physical screening protocol \cite{mattergen}. The accompanying source audit records the additional approaches examined, their conditioning information, their treatment of geometry, and which stages would have to be charged in a compatible comparison.

\section{Wyckoff charts, equivariance and evaluation}\label{app:wyckoff}
\paragraph{Setting and orbit decoder.}
For a conventional cell we use one specified Hall setting per group, tabulated by PyXtal 1.0.9 \cite{pyxtal}. If the row coordinate of an orbit representative is $s(u)=a+uB$, its affine chart has $\operatorname{rank}B\in\{0,1,2,3\}$. Coordinates are reduced modulo one. The table supplies the maps $s_j(u)=a_j+u B_j$ produced by acting on $s$ with one representative of each coset of its site stabilizer. For a right action, $sg=sh$ holds precisely when $gh^{-1}$ is in the stabilizer. Each independent occupied orbit is drawn once, with a common chemical element on all images. We reject coincident orbit sites or collisions between occupied orbits to fractional tolerance $10^{-7}$ and redraw the parameters of the new orbit up to 32 times. The program reports a failed proposal if this bound is reached. Its random continuous prior gives exceptional boundaries zero probability in the free directions. Fixed zero-degree charts and conventional-cell centering translations are included explicitly. A specialization or additional symmetry can still occur in the generated geometry.

For a row point operation $R$, compatibility requires $RGR^\top=G$. Given an SPD proposal $H$, the average $\overline G=|\mathcal G|^{-1}\sum_{R,t}RHR^\top$ remains SPD and is invariant because multiplication permutes the operation list. We rescale to the sampled volume and take a Cholesky factor. A species-preserving fractional operation of a compatible seed induces an orthogonal physical frame action. The paired equivariance theorem applied to the corresponding site permutation and its fractional translation shows that its operations remain a subgroup of the output stabilizer. This implication includes inversion and improper rotations. It does not imply that the detected group is exactly the requested group, and an arbitrary continuous unconditioned seed has no prescribed nontrivial stabilizer.

\paragraph{Label extraction and training.}
We run spglib 2.5.0 at 0.01\,\AA\ on the actual MP-20 input structure, imposing the PyXtal Hall setting on a second call where necessary \cite{spglib}. The native and conventional atom maps are joined through their primitive site indices. We recover a representative of each species-labelled crystallographic orbit by solving its affine chart over the possible wraps and verify the expanded coset set with a minimum-distance assignment. Verification includes conventional-to-native species ratios, volume per site, the compatible metric, and the public primitive-cell StructureMatcher. We retain the material ID and the reason for every rejection.
The train, validation and test partitions retained 27,112/27,136, 9,039/9,047 and 9,039/9,046 round trips, respectively. All 39 rejections failed the final StructureMatcher check. All 128 nominated test targets passed. Changing the symmetry-finding tolerance from 0.01 to 0.001 and 0.1\,\AA\ changed the group type on three and zero nominated test records, respectively. The corresponding validation counts were six and five. The release gives orbit counts by crystal system and free-coordinate dimension.
The common metric prior samples the shape of a training cell with the same group and Wyckoff-letter multiset when available, falls back first to group and then to the training archive, and applies the group average. A ridge regression on training element fractions and log site count models log volume per atom, with an empirical donor residual. Uniform free parameters supply the site coordinates. The conditional angular objective samples four training rows with replacement per update. For each row's exact labelled template it independently draws three data references, three generated negatives and three live queries. The positive-sampling archive contains 26,821 distinct templates, of which 26,591 have only one record. The supervised template establishes orbit multiplicity and site dimension; repeated sampling in a singleton bank supplies no additional geometries. Both C objectives use the same 6,723-parameter physical generator, shared train-only scaling, AdamW at $10^{-3}$, the same initialization and training draws, and 512 updates. We select among updates 0, 64, 256 and 512 by validation MR at five draws on 64 fixed IDs, breaking ties by lower matched RMS. The global and orbit-augmented objectives selected 256 and 512 updates, respectively.

\paragraph{Orbit statistic and independent checks.}
For each center and order $\ell=0,1,2$ we sum the real harmonics of the weighted neighbor vectors and contract over their components. We average this scalar separately over each occupied orbit. Group orbits by their element and their multiplicity. Within each class compute the variance of these means, then average over the occupied orbits and concatenate the three standardized entries with the physical fingerprint. A simultaneous transport of sites and orbit membership preserves each per-orbit mean. Reordering the orbits permutes terms inside the symmetric reduction. A change of setting transports these same orbit sets and preserves their multiplicities under a unimodular transformation. This statistic avoids an ordering by origin-dependent Wyckoff letters. Cell enlargement or merging occupations changes its domain. A categorical template label is constant within a template bank. The extra block supplies continuous information whenever equivalent-size orbits of the same species have distinct environments.
All 1,731 tabulated charts in 230 settings reproduced the orbit of a generic point under independently supplied spglib Hall operations. Maximum torus-coordinate discrepancy was $3.4\times10^{-16}$. Thirty-nine hand-chosen and random templates tested fixed inversion centers, mirror planes, free lines, a screw, a glide, centering and repeated same-species orbits under a restoring checkpoint, a forward call, rebasing and an optimizer update. The largest paired cell or site difference was $3.6\times10^{-15}$ and the added statistics agreed within $2.1\times10^{-16}$.
\paragraph{Detection and relaxation.}
For each tolerance we compare every requested operation with the independently detected spglib operations in the supplied setting, pairing equal fractional point matrices and requiring Cartesian translation discrepancies below the tolerance. We also compare the requested and detected international group numbers. As a separate direct check we expand each orbit from its first generated site, compare the two finite sets in both directions and test the isometry defect of each cell action. For $Q=L^{-1}RL$, we use $\max_j|\sigma_j(Q)-1|$ times the longest supplied lattice vector as a metric residual in \AA. The motif residual uses Cartesian distances at the assigned fractional wraps, and the check uses the larger residual. On every audited raw and relaxed cell, $0.1\max_j\|(L^{-1})_{:j}\|<0.044$, so every image vector shorter than a tested tolerance lies inside these wraps. For resolved occupancy we also require the shortest intersite or nonzero image distance to exceed the tolerance, so sites that coalesce there cannot count as separate occupants. We transport the initial origin by the mean wrapped site displacement when analysing the unconstrained final cell. All candidates retain their original indices. The explicit parameterized subgroup and the numerical recovery by symmetry finding can therefore be checked independently. With CHGNet 0.3.0 we relax sites and cell at zero applied pressure, using 100 FIRE steps and a generalized force threshold of 0.1 eV/\AA. Screened groups and raw groups have different sample sizes. No competing phase energies enter this retention experiment.

\begin{figure}[!htbp]\centering\includegraphics[width=.76\linewidth]{wyckoff-example.pdf}
\caption{The first multispecies nontrivial-group C-angular draw in the preselected relaxation cohort, for the template of MP-20 record mp-20959 (group 127). The two views share a physical scale. Raw and optimized coordinates and forces are archived. The raw nearest distance was 2.35\,\AA\ and the final distance was 2.63\,\AA.}\label{fig:wyckoff-example}
\end{figure}
\begin{table}[!htbp]\centering\footnotesize
\caption{C-trained angular predictions under three tolerances. Counts refer to all 2,560 raw draws from 128 templates, and 16 final CHGNet geometries from the first draw of the first 16 templates. The orbit column checks the expected occupancy within each supplied orbit under its own requested operations, including repeated occupations.}\label{tab:wyckoff-tolerance}
\begin{tabular}{lrrrrrrrr}\hline
 & \multicolumn{4}{c}{Raw (2,560)} & \multicolumn{4}{c}{After relaxation (16)}\\
Tolerance (\AA)&Subgroup&Set&Type&Orbit&Subgroup&Set&Type&Orbit\\\hline
.001 & 2,560 & 2,560 & 2,560 & 2,559 & 14 & 14 & 14 & 14\\
.01 & 2,557 & 2,556 & 2,556 & 2,553 & 15 & 15 & 15 & 15\\
.1 & 2,527 & 2,523 & 2,522 & 2,511 & 15 & 15 & 15 & 15\\
\hline\end{tabular}\end{table}
We match requested translations injectively to detected translations for each equal rotation. Containment together with equal operation counts, including centering translations, gives equality of these finite sets modulo lattice translations at the stated tolerance. At 0.001, 0.01 and 0.1\,\AA, finite-set equality and international type equality differed on 0, 0 and 6 of 12,928 raw candidates. Their corresponding disagreement counts after relaxation were 0, 0 and 0 of 96.
The external-comparator cohort contained 21 tetragonal, 30 orthorhombic, 33 cubic, 10 hexagonal, 13 trigonal, 17 monoclinic, 4 triclinic targets. The occupied orbits included 206 with 0, 162 with 1, 109 with 2, 110 with 3 free coordinates. The archived stratification counts every template under each degree it contains.
Independent detection missed the requested subgroup at 0.01\,\AA\ on 2 tetragonal draws and 1 orthorhombic draw.
On one Intel Core i7-12700H CPU thread, median batch-one latency across the first 16 test templates was 1.06 ms for the prior, 5.78 ms for C angular, and 4127.6 ms for the 40-call DiffCSP++ sampler. At batch 16, their measured throughputs were 948.7, 483.0 and 0.2 cells/s, respectively. The clock covers supplied-template decoding and all model calls through returned host arrays after two warm-ups. Other CPU jobs shared the host. Independently, generating the archived 2,560 respaced DiffCSP++ candidates took 187.9 s on the RTX 3060 Laptop GPU. The 128 native-schedule draws took 399.8 s.

\section{What the earlier Fourier fingerprint can recover}
For a fixed basis let $a_s(h)$ denote the average of $\exp(2\pi\mathrm{i}f_jh)$ over sites of element $s$. In a two-site AB crystal the cross-power at each coordinate vector $e_i$ is $\overline{a_A(e_i)}a_B(e_i)=\exp(2\pi\mathrm{i}(f_B-f_A)_i)$. Its phase gives the relative motif modulo the lattice. The six invertible log-metric entries give the cell up to its orthogonal frame. Under unimodular rebasing an individual frequency must transform as $h'=Uh$. A finite list of frequencies indexed by the deposited basis is generally changed.

For an explicit failure at eight sites, use the classical homometric subsets $S=\{0,1,4,6\}$ and $T=\{0,1,3,7\}$ of $\mathbb Z/12$ \cite{zrelation,homometry}. For each set place four A atoms at $(s/12,0,0)$ and four B atoms at a common translate $(0,1/3,0)$ of this row. Both sets have four ordered differences at zero, two at six and one at every other residue. Hence
\[
 \left|a_A(h)\right|^2=\tfrac1{16}\sum_{s,t}\exp(2\pi\mathrm{i}(s-t)h_1/12)
\]
agrees at every integral wave. Each B amplitude is the A amplitude times the same known phase in both crystals. In a trained mixed triple at $(e_i,e_j,e_i+e_j)$ with $i\ne j$, at most one factor from the first two frequencies probes the A row. The motif-dependent factor reduces to unity or a power. In the multielement model, the whole-motif and first-species reference amplitudes also reduce to the A amplitude times a common factor. These observations give equality of all 70 binary and all 1,306 multielement entries, including the common cell. The cyclic gap multisets in the A rows are $\{1,2,3,6\}$ and $\{1,2,4,5\}$. They exclude translation or reversal. Taking a rectangular translation lattice with rationally independent squared axis lengths forces every lattice isometry to map each axis to itself up to sign. Both decorations have the displayed full translation periods. The examples are therefore physically inequivalent.

The real part of $a_A(e_1)^2\overline{a_B(2e_1)}/\sqrt6$ restores a piece of the missing phase. Substituting the four twelfth roots of unity of each set gives an absolute difference $1/(16\sqrt2)=0.044194\ldots$. The baseline independently evaluated the implemented blocks and the extra channel in double precision. The physical powers of this study use the correlations of edges meeting at each atom. Figure~\ref{fig:v2-witness} tests this different construction, and the eight-site results in Section~\ref{sec:v2checks} use the very same coordinates and decorations as the earlier Fourier test.

\paragraph{Relation to the broader benchmark.}
The fixed-basis release compares an anchor retained in the first positive slot against all-independent positive sampling on the same composition schedule. It trains the compact drift, independent-bank drift, flow and dedicated MeanFlow for one, three and ten passes at seeds 11, 23 and 37. Its WBM experiments use formula-disjoint splits, and its additional Alexandria allocation test excludes earlier IDs and reduced formulas. The CPU binary study has its own 128 targets. The public 9,046-structure MP-20 split was already inspected in that work. These prior training, volume correction and transfer results apply to their original fixed-basis checkpoints. Lang2Str adds inferred textual and symmetry information before geometric decoding; the published oracle variant gives a different input \cite{lang2str}.

\section{Sensitivity to maximum site displacement and native volume}\label{app:csp-sensitivity}
We rescore the ordered raw proposals under two stricter criteria on the same 9,046 official composition-conditioned targets. The maximum-site criterion uses the public \texttt{StructureMatcher.fit} acceptance under the prescribed primitive-cell reduction, volume scaling, exact species, lattice tolerance 0.3, site tolerance 0.5 and angle tolerance $10^\circ$. It requires an accepted mapping whose largest normalized displacement passes the site tolerance. The volume-qualified criterion requires the standard RMS match and $|\log(V_{\rm draw}/V_{\rm ref})|\leq\log(1.10)$ on the same draw, using the original cell volumes before matcher scaling. Their compositions, and hence site counts, are identical.
\begin{table}[!htbp]\centering\scriptsize
\caption{Raw recovery (\%) under stricter structure criteria. Differences are percentage points, with 95\% paired target-bootstrap intervals. Each target carries the same three training or prior seeds 11, 23 and 37 and up to twenty ordered draws.}\label{tab:csp-sensitivity}
\setlength{\tabcolsep}{3.3pt}\begin{tabular}{lrrrrrr}\hline
Criterion & $k$ & Prior & Radial & Angular & Angular $-$ prior & Angular $-$ radial\\\hline
Maximum site & 1 & 1.86 & 2.58 & 2.68 & +0.82 [+0.70,+0.94] & +0.10 [+0.05,+0.15]\\
 & 5 & 6.22 & 8.74 & 9.05 & +2.83 [+2.59,+3.07] & +0.31 [+0.22,+0.40]\\
 & 20 & 13.47 & 17.39 & 17.81 & +4.34 [+4.05,+4.63] & +0.42 [+0.32,+0.52]\\
Native volume & 1 & 2.61 & 3.28 & 2.79 & +0.18 [-0.02,+0.40] & -0.49 [-0.70,-0.27]\\
 & 5 & 10.02 & 12.28 & 10.06 & +0.04 [-0.31,+0.41] & -2.22 [-2.56,-1.86]\\
 & 20 & 21.15 & 24.60 & 21.24 & +0.08 [-0.31,+0.46] & -3.37 [-3.74,-2.98]\\
\hline\end{tabular}\end{table}
Across twenty raw draws per target and seed, the mean absolute log-volume errors for the prior, radial and angular maps were 0.116, 0.120, 0.164, respectively. The corresponding mean shape errors were 0.435, 0.435, 0.435. Shape error is the Euclidean distance between sorted eigenvalues of $\frac12\log(LL^\top)$ after subtraction of their mean, computed in each deposited basis. This removes uniform scale and Cartesian orientation. A general unimodular rebasing can change this diagnostic.
On the first draw of every target and seed, the mean fractional motif RMS errors in the same order were 0.308, 0.297, 0.296. For this diagnostic we use the deposited basis, align a center of the rarest species, alternate same-species assignments with a common origin refinement for up to four rounds, and retain the best feasible displacement. Each displacement is wrapped componentwise. This diagnostic is dimensionless and depends on the basis and alignment search. It includes targets with no matcher hit. The archive records the individual-seed results, absolute intervals and paired intervals for all three geometry diagnostics.

\end{document}
